% --- Archon physics formalization source begin ---
% archon:physics
% archon:covers PhyXMiniProblems/problem_phyx_mini_0898.lean
% archon:source-report reports/phyx_mini/problem_phyx_mini_0898.source.json

\chapter{Physics problem phyx\_mini\_0898}
\label{ch:PhyXMiniProblems_problem_phyx_mini_0898}

\paragraph{Problem source.}
\#\# Question

What is the magnitude of the force \textbackslash{}( \textbackslash{}vec\{F\} \textbackslash{}) on the \textbackslash{}( 5.0\textbackslash{},\textbackslash{}mathrm\{nC\} \textbackslash{}) charge in figure?

\#\# Answer choices

- A: A: 7.25 \textbackslash{}times 10\textasciicircum{}\{-4\}N

- B: B: 6.35 \textbackslash{}times 10\textasciicircum{}\{-4\}N

- C: C: 1.7 \textbackslash{}times 10\textasciicircum{}\{-4\}N

- D: D: 1.33 \textbackslash{}times 10\textasciicircum{}\{-4\}N

\#\# Figure caption (auxiliary; use the image as primary evidence)

The image depicts a diagram with three charged particles arranged in a right triangle formation. Here's a detailed description:

- **Top Left:** A red circle with a plus sign inside representing a positive charge labeled "10 nC."
- **Top Right:** A blue circle with a minus sign inside representing a negative charge labeled "-10 nC."
- **Bottom Left:** A red circle with a plus sign inside representing a positive charge labeled "5.0 nC."
- **Connection Lines:** 
  - A horizontal dashed line connects the top two charges, labeled as "3.0 cm."
  - A vertical dashed line connects the top left and bottom left charges, labeled as "4.0 cm."
  - A dashed line connects the top right and bottom left charges, forming the hypotenuse of the triangle.

Each charge is enclosed within its respective colored circle, indicating their nature (positive or negative), and lines denote the spatial relationships and distances between them.

\#\# Dataset metadata

Electrostatics; Spatial Relation Reasoning, Multi-Formula Reasoning

\paragraph{Recorded answer/context.}
C: C: 1.7 \textbackslash{}times 10\textasciicircum{}\{-4\}N

\paragraph{Figure/image path.}
/root/proposal\_for\_physic/science-mango/phyx\_mini\_run/phyx\_data/test\_image/898.png

\paragraph{Formalization target.}
create a compiling Lean file with sorry bodies at `PhyXMiniProblems/problem\_phyx\_mini\_0898.lean`. The Lean declarations must preserve the physical quantities, dimensions or dimensional roles, figure labels, governing-law hypotheses, and final relation expressed by this problem.
Use Mathlib/Physlib names found through LeanExplore where available. If a physics API is missing, introduce faithful local abstractions rather than scalar placeholder aliases.

\begin{theorem}[Physics formalization target]
\label{thm:physics:phyx_mini_0898:target}
\lean{PhyXMiniProblems.ProblemPhyXMini0898.problem_phyx_mini_0898}
\uses{def:physics:phyx-mini-0898:phyxminiproblems-problemphyxmini0898-forcemagnitudeinnewtons, def:physics:phyx-mini-0898:phyxminiproblems-problemphyxmini0898-threepointchargeforcesetup, def:physics:phyx-mini-0898:phyxminiproblems-problemphyxmini0898-matchessuppliedthreechargefigure, def:physics:phyx-mini-0898:phyxminiproblems-problemphyxmini0898-usesschoolcoulombconstant, def:physics:phyx-mini-0898:phyxminiproblems-problemphyxmini0898-hasphysicalthreechargeparameters, def:physics:phyx-mini-0898:phyxminiproblems-problemphyxmini0898-satisfiespairwisecoulombforcelaw, def:physics:phyx-mini-0898:phyxminiproblems-problemphyxmini0898-satisfieselectrostaticforcesuperposition, def:physics:phyx-mini-0898:phyxminiproblems-problemphyxmini0898-recordeddatasetanswer, def:physics:phyx-mini-0898:phyxminiproblems-problemphyxmini0898-roundstonearesttenmicronewtons, def:physics:phyx-mini-0898:phyxminiproblems-problemphyxmini0898-isuniquenearestdisplayedforce}
The assigned autoformalize agent should translate this physics problem into Lean declarations in the covered file, with theorem and lemma proof bodies written as `by sorry`.
\end{theorem}
\begin{proof}
This is an autoformalization task, not a proof task. Produce statements that can later be proved by the physics prover without weakening the physical model.
\end{proof}

% --- Archon declaration topology begin ---
\section{Declaration topology}
\begin{definition}[Planar Vector]
  \label{def:physics:phyx-mini-0898:phyxminiproblems-problemphyxmini0898-planarvector}
  \lean{PhyXMiniProblems.ProblemPhyXMini0898.PlanarVector}
  A Cartesian vector in the plane of the supplied figure
\end{definition}
\begin{proof}
  This is the declared model definition of Planar Vector.
\end{proof}
\begin{definition}[force Dimension]
  \label{def:physics:phyx-mini-0898:phyxminiproblems-problemphyxmini0898-forcedimension}
  \lean{PhyXMiniProblems.ProblemPhyXMini0898.forceDimension}
  The physical dimension M L T⁻² of force
\end{definition}
\begin{proof}
  This is the declared model definition of force Dimension.
\end{proof}
\begin{definition}[Length Quantity]
  \label{def:physics:phyx-mini-0898:phyxminiproblems-problemphyxmini0898-lengthquantity}
  \lean{PhyXMiniProblems.ProblemPhyXMini0898.LengthQuantity}
  A nonnegative, unit-independent physical length
\end{definition}
\begin{proof}
  This is the declared model definition of Length Quantity.
\end{proof}
\begin{definition}[Signed Charge Quantity]
  \label{def:physics:phyx-mini-0898:phyxminiproblems-problemphyxmini0898-signedchargequantity}
  \lean{PhyXMiniProblems.ProblemPhyXMini0898.SignedChargeQuantity}
  A signed, unit-independent electric charge
\end{definition}
\begin{proof}
  This is the declared model definition of Signed Charge Quantity.
\end{proof}
\begin{definition}[Planar Position Quantity]
  \label{def:physics:phyx-mini-0898:phyxminiproblems-problemphyxmini0898-planarpositionquantity}
  \lean{PhyXMiniProblems.ProblemPhyXMini0898.PlanarPositionQuantity}
  \uses{def:physics:phyx-mini-0898:phyxminiproblems-problemphyxmini0898-planarvector}
  A unit-independent physical position in the plane
\end{definition}
\begin{proof}
  This is the declared model definition of Planar Position Quantity.
\end{proof}
\begin{definition}[Planar Force Quantity]
  \label{def:physics:phyx-mini-0898:phyxminiproblems-problemphyxmini0898-planarforcequantity}
  \lean{PhyXMiniProblems.ProblemPhyXMini0898.PlanarForceQuantity}
  \uses{def:physics:phyx-mini-0898:phyxminiproblems-problemphyxmini0898-planarvector, def:physics:phyx-mini-0898:phyxminiproblems-problemphyxmini0898-forcedimension}
  A unit-independent physical force vector in the plane
\end{definition}
\begin{proof}
  This is the declared model definition of Planar Force Quantity.
\end{proof}
\begin{definition}[length In Meters]
  \label{def:physics:phyx-mini-0898:phyxminiproblems-problemphyxmini0898-lengthinmeters}
  \lean{PhyXMiniProblems.ProblemPhyXMini0898.lengthInMeters}
  \uses{def:physics:phyx-mini-0898:phyxminiproblems-problemphyxmini0898-lengthquantity}
  Coherent-SI metre readout of a nonnegative physical length
\end{definition}
\begin{proof}
  This is the declared model definition of length In Meters.
\end{proof}
\begin{definition}[length In Centimeters]
  \label{def:physics:phyx-mini-0898:phyxminiproblems-problemphyxmini0898-lengthincentimeters}
  \lean{PhyXMiniProblems.ProblemPhyXMini0898.lengthInCentimeters}
  \uses{def:physics:phyx-mini-0898:phyxminiproblems-problemphyxmini0898-lengthquantity, def:physics:phyx-mini-0898:phyxminiproblems-problemphyxmini0898-lengthinmeters}
  Centimetre readout used by the two dimension labels in the figure
\end{definition}
\begin{proof}
  This is the declared model definition of length In Centimeters.
\end{proof}
\begin{definition}[charge In Coulombs]
  \label{def:physics:phyx-mini-0898:phyxminiproblems-problemphyxmini0898-chargeincoulombs}
  \lean{PhyXMiniProblems.ProblemPhyXMini0898.chargeInCoulombs}
  \uses{def:physics:phyx-mini-0898:phyxminiproblems-problemphyxmini0898-signedchargequantity}
  Coherent-SI coulomb readout of a signed physical charge
\end{definition}
\begin{proof}
  This is the declared model definition of charge In Coulombs.
\end{proof}
\begin{definition}[charge In Nanocoulombs]
  \label{def:physics:phyx-mini-0898:phyxminiproblems-problemphyxmini0898-chargeinnanocoulombs}
  \lean{PhyXMiniProblems.ProblemPhyXMini0898.chargeInNanocoulombs}
  \uses{def:physics:phyx-mini-0898:phyxminiproblems-problemphyxmini0898-signedchargequantity, def:physics:phyx-mini-0898:phyxminiproblems-problemphyxmini0898-chargeincoulombs}
  Nanocoulomb readout used by the three printed charge labels
\end{definition}
\begin{proof}
  This is the declared model definition of charge In Nanocoulombs.
\end{proof}
\begin{definition}[position Vector In Meters]
  \label{def:physics:phyx-mini-0898:phyxminiproblems-problemphyxmini0898-positionvectorinmeters}
  \lean{PhyXMiniProblems.ProblemPhyXMini0898.positionVectorInMeters}
  \uses{def:physics:phyx-mini-0898:phyxminiproblems-problemphyxmini0898-planarvector, def:physics:phyx-mini-0898:phyxminiproblems-problemphyxmini0898-planarpositionquantity}
  Cartesian metre readout of a physical planar position
\end{definition}
\begin{proof}
  This is the declared model definition of position Vector In Meters.
\end{proof}
\begin{definition}[force Vector In Newtons]
  \label{def:physics:phyx-mini-0898:phyxminiproblems-problemphyxmini0898-forcevectorinnewtons}
  \lean{PhyXMiniProblems.ProblemPhyXMini0898.forceVectorInNewtons}
  \uses{def:physics:phyx-mini-0898:phyxminiproblems-problemphyxmini0898-planarvector, def:physics:phyx-mini-0898:phyxminiproblems-problemphyxmini0898-planarforcequantity}
  Cartesian newton readout of a physical planar force
\end{definition}
\begin{proof}
  This is the declared model definition of force Vector In Newtons.
\end{proof}
\begin{definition}[force Magnitude In Newtons]
  \label{def:physics:phyx-mini-0898:phyxminiproblems-problemphyxmini0898-forcemagnitudeinnewtons}
  \lean{PhyXMiniProblems.ProblemPhyXMini0898.forceMagnitudeInNewtons}
  \uses{def:physics:phyx-mini-0898:phyxminiproblems-problemphyxmini0898-planarforcequantity, def:physics:phyx-mini-0898:phyxminiproblems-problemphyxmini0898-forcevectorinnewtons}
  Euclidean magnitude, in newtons, of a physical planar force
\end{definition}
\begin{proof}
  This is the declared model definition of force Magnitude In Newtons.
\end{proof}
\begin{definition}[Rectangle Corner]
  \label{def:physics:phyx-mini-0898:phyxminiproblems-problemphyxmini0898-rectanglecorner}
  \lean{PhyXMiniProblems.ProblemPhyXMini0898.RectangleCorner}
  The four corners of the dashed rectangle in image 898.png
\end{definition}
\begin{proof}
  This is the declared model definition of Rectangle Corner.
\end{proof}
\begin{definition}[Charge Label]
  \label{def:physics:phyx-mini-0898:phyxminiproblems-problemphyxmini0898-chargelabel}
  \lean{PhyXMiniProblems.ProblemPhyXMini0898.ChargeLabel}
  The three physical charges, named by value and role in the figure
\end{definition}
\begin{proof}
  This is the declared model definition of Charge Label.
\end{proof}
\begin{definition}[Source Charge]
  \label{def:physics:phyx-mini-0898:phyxminiproblems-problemphyxmini0898-sourcecharge}
  \lean{PhyXMiniProblems.ProblemPhyXMini0898.SourceCharge}
  The two charges exerting force on the lower-left target charge
\end{definition}
\begin{proof}
  This is the declared model definition of Source Charge.
\end{proof}
\begin{definition}[source Charge Label]
  \label{def:physics:phyx-mini-0898:phyxminiproblems-problemphyxmini0898-sourcechargelabel}
  \lean{PhyXMiniProblems.ProblemPhyXMini0898.sourceChargeLabel}
  \uses{def:physics:phyx-mini-0898:phyxminiproblems-problemphyxmini0898-chargelabel, def:physics:phyx-mini-0898:phyxminiproblems-problemphyxmini0898-sourcecharge}
  Charge label associated with each source role
\end{definition}
\begin{proof}
  This is the declared model definition of source Charge Label.
\end{proof}
\begin{definition}[target Charge Label]
  \label{def:physics:phyx-mini-0898:phyxminiproblems-problemphyxmini0898-targetchargelabel}
  \lean{PhyXMiniProblems.ProblemPhyXMini0898.targetChargeLabel}
  \uses{def:physics:phyx-mini-0898:phyxminiproblems-problemphyxmini0898-chargelabel}
  The charge whose force is requested
\end{definition}
\begin{proof}
  This is the declared model definition of target Charge Label.
\end{proof}
\begin{definition}[expected Charge Corner]
  \label{def:physics:phyx-mini-0898:phyxminiproblems-problemphyxmini0898-expectedchargecorner}
  \lean{PhyXMiniProblems.ProblemPhyXMini0898.expectedChargeCorner}
  \uses{def:physics:phyx-mini-0898:phyxminiproblems-problemphyxmini0898-rectanglecorner, def:physics:phyx-mini-0898:phyxminiproblems-problemphyxmini0898-chargelabel}
  Corner occupied by each charge in the primary image
\end{definition}
\begin{proof}
  This is the declared model definition of expected Charge Corner.
\end{proof}
\begin{definition}[Figure Charge Sign]
  \label{def:physics:phyx-mini-0898:phyxminiproblems-problemphyxmini0898-figurechargesign}
  \lean{PhyXMiniProblems.ProblemPhyXMini0898.FigureChargeSign}
  Sign glyph drawn inside a charge circle
\end{definition}
\begin{proof}
  This is the declared model definition of Figure Charge Sign.
\end{proof}
\begin{definition}[Figure Charge Color]
  \label{def:physics:phyx-mini-0898:phyxminiproblems-problemphyxmini0898-figurechargecolor}
  \lean{PhyXMiniProblems.ProblemPhyXMini0898.FigureChargeColor}
  Fill colour used to distinguish positive and negative charges
\end{definition}
\begin{proof}
  This is the declared model definition of Figure Charge Color.
\end{proof}
\begin{definition}[expected Charge Sign]
  \label{def:physics:phyx-mini-0898:phyxminiproblems-problemphyxmini0898-expectedchargesign}
  \lean{PhyXMiniProblems.ProblemPhyXMini0898.expectedChargeSign}
  \uses{def:physics:phyx-mini-0898:phyxminiproblems-problemphyxmini0898-chargelabel, def:physics:phyx-mini-0898:phyxminiproblems-problemphyxmini0898-figurechargesign}
  Sign glyph expected for each of the three charges
\end{definition}
\begin{proof}
  This is the declared model definition of expected Charge Sign.
\end{proof}
\begin{definition}[expected Charge Color]
  \label{def:physics:phyx-mini-0898:phyxminiproblems-problemphyxmini0898-expectedchargecolor}
  \lean{PhyXMiniProblems.ProblemPhyXMini0898.expectedChargeColor}
  \uses{def:physics:phyx-mini-0898:phyxminiproblems-problemphyxmini0898-chargelabel, def:physics:phyx-mini-0898:phyxminiproblems-problemphyxmini0898-figurechargecolor}
  Circle colour expected for each of the three charges
\end{definition}
\begin{proof}
  This is the declared model definition of expected Charge Color.
\end{proof}
\begin{definition}[expected Charge In Nanocoulombs]
  \label{def:physics:phyx-mini-0898:phyxminiproblems-problemphyxmini0898-expectedchargeinnanocoulombs}
  \lean{PhyXMiniProblems.ProblemPhyXMini0898.expectedChargeInNanocoulombs}
  \uses{def:physics:phyx-mini-0898:phyxminiproblems-problemphyxmini0898-chargelabel}
  Signed nanocoulomb number printed beside each charge
\end{definition}
\begin{proof}
  This is the declared model definition of expected Charge In Nanocoulombs.
\end{proof}
\begin{definition}[Rectangle Edge]
  \label{def:physics:phyx-mini-0898:phyxminiproblems-problemphyxmini0898-rectangleedge}
  \lean{PhyXMiniProblems.ProblemPhyXMini0898.RectangleEdge}
  The four dashed sides of the rectangle
\end{definition}
\begin{proof}
  This is the declared model definition of Rectangle Edge.
\end{proof}
\begin{definition}[Dimension Label Side]
  \label{def:physics:phyx-mini-0898:phyxminiproblems-problemphyxmini0898-dimensionlabelside}
  \lean{PhyXMiniProblems.ProblemPhyXMini0898.DimensionLabelSide}
  The two sides carrying numerical dimension labels
\end{definition}
\begin{proof}
  This is the declared model definition of Dimension Label Side.
\end{proof}
\begin{definition}[Three Charge Rectangle Figure]
  \label{def:physics:phyx-mini-0898:phyxminiproblems-problemphyxmini0898-threechargerectanglefigure}
  \lean{PhyXMiniProblems.ProblemPhyXMini0898.ThreeChargeRectangleFigure}
  \uses{def:physics:phyx-mini-0898:phyxminiproblems-problemphyxmini0898-rectanglecorner, def:physics:phyx-mini-0898:phyxminiproblems-problemphyxmini0898-chargelabel, def:physics:phyx-mini-0898:phyxminiproblems-problemphyxmini0898-figurechargesign, def:physics:phyx-mini-0898:phyxminiproblems-problemphyxmini0898-figurechargecolor, def:physics:phyx-mini-0898:phyxminiproblems-problemphyxmini0898-rectangleedge, def:physics:phyx-mini-0898:phyxminiproblems-problemphyxmini0898-dimensionlabelside}
  Literal presentation data transcribed from the primary bitmap. It contains no force readout and no answer-choice label
\end{definition}
\begin{proof}
  This is the declared model definition of Three Charge Rectangle Figure.
\end{proof}
\begin{definition}[Three Point Charge Force Setup]
  \label{def:physics:phyx-mini-0898:phyxminiproblems-problemphyxmini0898-threepointchargeforcesetup}
  \lean{PhyXMiniProblems.ProblemPhyXMini0898.ThreePointChargeForceSetup}
  \uses{def:physics:phyx-mini-0898:phyxminiproblems-problemphyxmini0898-lengthquantity, def:physics:phyx-mini-0898:phyxminiproblems-problemphyxmini0898-signedchargequantity, def:physics:phyx-mini-0898:phyxminiproblems-problemphyxmini0898-planarpositionquantity, def:physics:phyx-mini-0898:phyxminiproblems-problemphyxmini0898-planarforcequantity, def:physics:phyx-mini-0898:phyxminiproblems-problemphyxmini0898-rectanglecorner, def:physics:phyx-mini-0898:phyxminiproblems-problemphyxmini0898-chargelabel, def:physics:phyx-mini-0898:phyxminiproblems-problemphyxmini0898-sourcecharge, def:physics:phyx-mini-0898:phyxminiproblems-problemphyxmini0898-threechargerectanglefigure}
  Independent physical quantities and force observables. The resultant force is not defined from the recorded answer; the governing-law premises below relate it to the two independent pair forces
\end{definition}
\begin{proof}
  This is the declared model definition of Three Point Charge Force Setup.
\end{proof}
\begin{definition}[charge Position In Meters]
  \label{def:physics:phyx-mini-0898:phyxminiproblems-problemphyxmini0898-chargepositioninmeters}
  \lean{PhyXMiniProblems.ProblemPhyXMini0898.chargePositionInMeters}
  \uses{def:physics:phyx-mini-0898:phyxminiproblems-problemphyxmini0898-planarvector, def:physics:phyx-mini-0898:phyxminiproblems-problemphyxmini0898-positionvectorinmeters, def:physics:phyx-mini-0898:phyxminiproblems-problemphyxmini0898-chargelabel, def:physics:phyx-mini-0898:phyxminiproblems-problemphyxmini0898-threepointchargeforcesetup}
  Position of a named charge in coherent-SI metres
\end{definition}
\begin{proof}
  This is the declared model definition of charge Position In Meters.
\end{proof}
\begin{definition}[source Position In Meters]
  \label{def:physics:phyx-mini-0898:phyxminiproblems-problemphyxmini0898-sourcepositioninmeters}
  \lean{PhyXMiniProblems.ProblemPhyXMini0898.sourcePositionInMeters}
  \uses{def:physics:phyx-mini-0898:phyxminiproblems-problemphyxmini0898-planarvector, def:physics:phyx-mini-0898:phyxminiproblems-problemphyxmini0898-sourcecharge, def:physics:phyx-mini-0898:phyxminiproblems-problemphyxmini0898-sourcechargelabel, def:physics:phyx-mini-0898:phyxminiproblems-problemphyxmini0898-threepointchargeforcesetup, def:physics:phyx-mini-0898:phyxminiproblems-problemphyxmini0898-chargepositioninmeters}
  Position of a named source charge in coherent-SI metres
\end{definition}
\begin{proof}
  This is the declared model definition of source Position In Meters.
\end{proof}
\begin{definition}[target Position In Meters]
  \label{def:physics:phyx-mini-0898:phyxminiproblems-problemphyxmini0898-targetpositioninmeters}
  \lean{PhyXMiniProblems.ProblemPhyXMini0898.targetPositionInMeters}
  \uses{def:physics:phyx-mini-0898:phyxminiproblems-problemphyxmini0898-planarvector, def:physics:phyx-mini-0898:phyxminiproblems-problemphyxmini0898-targetchargelabel, def:physics:phyx-mini-0898:phyxminiproblems-problemphyxmini0898-threepointchargeforcesetup, def:physics:phyx-mini-0898:phyxminiproblems-problemphyxmini0898-chargepositioninmeters}
  Position of the lower-left target charge in coherent-SI metres
\end{definition}
\begin{proof}
  This is the declared model definition of target Position In Meters.
\end{proof}
\begin{definition}[source To Target Displacement In Meters]
  \label{def:physics:phyx-mini-0898:phyxminiproblems-problemphyxmini0898-sourcetotargetdisplacementinmeters}
  \lean{PhyXMiniProblems.ProblemPhyXMini0898.sourceToTargetDisplacementInMeters}
  \uses{def:physics:phyx-mini-0898:phyxminiproblems-problemphyxmini0898-planarvector, def:physics:phyx-mini-0898:phyxminiproblems-problemphyxmini0898-sourcecharge, def:physics:phyx-mini-0898:phyxminiproblems-problemphyxmini0898-threepointchargeforcesetup, def:physics:phyx-mini-0898:phyxminiproblems-problemphyxmini0898-sourcepositioninmeters, def:physics:phyx-mini-0898:phyxminiproblems-problemphyxmini0898-targetpositioninmeters}
  Displacement from a source to the target. This orientation makes k q\_target q\_source r / ‖r‖³ point away for like charges and toward for unlike charges
\end{definition}
\begin{proof}
  This is the declared model definition of source To Target Displacement In Meters.
\end{proof}
\begin{definition}[Matches Supplied Three Charge Figure]
  \label{def:physics:phyx-mini-0898:phyxminiproblems-problemphyxmini0898-matchessuppliedthreechargefigure}
  \lean{PhyXMiniProblems.ProblemPhyXMini0898.MatchesSuppliedThreeChargeFigure}
  \uses{def:physics:phyx-mini-0898:phyxminiproblems-problemphyxmini0898-lengthinmeters, def:physics:phyx-mini-0898:phyxminiproblems-problemphyxmini0898-lengthincentimeters, def:physics:phyx-mini-0898:phyxminiproblems-problemphyxmini0898-chargeinnanocoulombs, def:physics:phyx-mini-0898:phyxminiproblems-problemphyxmini0898-positionvectorinmeters, def:physics:phyx-mini-0898:phyxminiproblems-problemphyxmini0898-expectedchargecorner, def:physics:phyx-mini-0898:phyxminiproblems-problemphyxmini0898-expectedchargesign, def:physics:phyx-mini-0898:phyxminiproblems-problemphyxmini0898-expectedchargecolor, def:physics:phyx-mini-0898:phyxminiproblems-problemphyxmini0898-expectedchargeinnanocoulombs, def:physics:phyx-mini-0898:phyxminiproblems-problemphyxmini0898-threepointchargeforcesetup}
  Facts read from image 898.png, together with their association to the dimensionful charges, lengths, and positions. No resultant-force value appears here
\end{definition}
\begin{proof}
  This is the declared model definition of Matches Supplied Three Charge Figure.
\end{proof}
\begin{definition}[Uses School Coulomb Constant]
  \label{def:physics:phyx-mini-0898:phyxminiproblems-problemphyxmini0898-usesschoolcoulombconstant}
  \lean{PhyXMiniProblems.ProblemPhyXMini0898.UsesSchoolCoulombConstant}
  \uses{def:physics:phyx-mini-0898:phyxminiproblems-problemphyxmini0898-threepointchargeforcesetup}
  Rounded free-space calibration used in the textbook multiple-choice calculation. Physlib's scalar has coherent-SI units N m² / C² here
\end{definition}
\begin{proof}
  This is the declared model definition of Uses School Coulomb Constant.
\end{proof}
\begin{definition}[Has Physical Three Charge Parameters]
  \label{def:physics:phyx-mini-0898:phyxminiproblems-problemphyxmini0898-hasphysicalthreechargeparameters}
  \lean{PhyXMiniProblems.ProblemPhyXMini0898.HasPhysicalThreeChargeParameters}
  \uses{def:physics:phyx-mini-0898:phyxminiproblems-problemphyxmini0898-lengthinmeters, def:physics:phyx-mini-0898:phyxminiproblems-problemphyxmini0898-chargeincoulombs, def:physics:phyx-mini-0898:phyxminiproblems-problemphyxmini0898-threepointchargeforcesetup, def:physics:phyx-mini-0898:phyxminiproblems-problemphyxmini0898-sourcetotargetdisplacementinmeters}
  Nondegeneracy and positivity conditions selecting the physical branch
\end{definition}
\begin{proof}
  This is the declared model definition of Has Physical Three Charge Parameters.
\end{proof}
\begin{definition}[Satisfies Pairwise Coulomb Force Law]
  \label{def:physics:phyx-mini-0898:phyxminiproblems-problemphyxmini0898-satisfiespairwisecoulombforcelaw}
  \lean{PhyXMiniProblems.ProblemPhyXMini0898.SatisfiesPairwiseCoulombForceLaw}
  \uses{def:physics:phyx-mini-0898:phyxminiproblems-problemphyxmini0898-chargeincoulombs, def:physics:phyx-mini-0898:phyxminiproblems-problemphyxmini0898-forcevectorinnewtons, def:physics:phyx-mini-0898:phyxminiproblems-problemphyxmini0898-sourcechargelabel, def:physics:phyx-mini-0898:phyxminiproblems-problemphyxmini0898-targetchargelabel, def:physics:phyx-mini-0898:phyxminiproblems-problemphyxmini0898-threepointchargeforcesetup, def:physics:phyx-mini-0898:phyxminiproblems-problemphyxmini0898-sourcetotargetdisplacementinmeters}
  Vector Coulomb law for the force of each source on the target. For displacement r from source to target, this is F = k q\_target q\_source r / ‖r‖³. It is a governing law and contains no requested numerical magnitude
\end{definition}
\begin{proof}
  This is the declared model definition of Satisfies Pairwise Coulomb Force Law.
\end{proof}
\begin{definition}[Satisfies Electrostatic Force Superposition]
  \label{def:physics:phyx-mini-0898:phyxminiproblems-problemphyxmini0898-satisfieselectrostaticforcesuperposition}
  \lean{PhyXMiniProblems.ProblemPhyXMini0898.SatisfiesElectrostaticForceSuperposition}
  \uses{def:physics:phyx-mini-0898:phyxminiproblems-problemphyxmini0898-forcevectorinnewtons, def:physics:phyx-mini-0898:phyxminiproblems-problemphyxmini0898-sourcecharge, def:physics:phyx-mini-0898:phyxminiproblems-problemphyxmini0898-threepointchargeforcesetup}
  The resultant force on the target is the vector sum of both pair forces
\end{definition}
\begin{proof}
  This is the declared model definition of Satisfies Electrostatic Force Superposition.
\end{proof}
\begin{definition}[Answer Choice]
  \label{def:physics:phyx-mini-0898:phyxminiproblems-problemphyxmini0898-answerchoice}
  \lean{PhyXMiniProblems.ProblemPhyXMini0898.AnswerChoice}
  Labels attached to the four multiple-choice answers
\end{definition}
\begin{proof}
  This is the declared model definition of Answer Choice.
\end{proof}
\begin{definition}[force Magnitude In Newtons]
  \label{def:physics:phyx-mini-0898:phyxminiproblems-problemphyxmini0898-answerchoice-forcemagnitudeinnewtons}
  \lean{PhyXMiniProblems.ProblemPhyXMini0898.AnswerChoice.forceMagnitudeInNewtons}
  \uses{def:physics:phyx-mini-0898:phyxminiproblems-problemphyxmini0898-forcemagnitudeinnewtons, def:physics:phyx-mini-0898:phyxminiproblems-problemphyxmini0898-answerchoice}
  Force magnitude in newtons printed beside each answer label
\end{definition}
\begin{proof}
  This is the declared model definition of force Magnitude In Newtons.
\end{proof}
\begin{definition}[recorded Dataset Answer]
  \label{def:physics:phyx-mini-0898:phyxminiproblems-problemphyxmini0898-recordeddatasetanswer}
  \lean{PhyXMiniProblems.ProblemPhyXMini0898.recordedDatasetAnswer}
  \uses{def:physics:phyx-mini-0898:phyxminiproblems-problemphyxmini0898-answerchoice}
  Answer label recorded by the supplied dataset metadata
\end{definition}
\begin{proof}
  This is the declared model definition of recorded Dataset Answer.
\end{proof}
\begin{definition}[Rounds To Nearest Ten Micro Newtons]
  \label{def:physics:phyx-mini-0898:phyxminiproblems-problemphyxmini0898-roundstonearesttenmicronewtons}
  \lean{PhyXMiniProblems.ProblemPhyXMini0898.RoundsToNearestTenMicroNewtons}
  A magnitude rounds to a displayed reading at the nearest 10⁻⁵ N; the half-step tolerance is 5 * 10⁻⁶ N
\end{definition}
\begin{proof}
  This is the declared model definition of Rounds To Nearest Ten Micro Newtons.
\end{proof}
\begin{definition}[Is Unique Nearest Displayed Force]
  \label{def:physics:phyx-mini-0898:phyxminiproblems-problemphyxmini0898-isuniquenearestdisplayedforce}
  \lean{PhyXMiniProblems.ProblemPhyXMini0898.IsUniqueNearestDisplayedForce}
  \uses{def:physics:phyx-mini-0898:phyxminiproblems-problemphyxmini0898-forcemagnitudeinnewtons, def:physics:phyx-mini-0898:phyxminiproblems-problemphyxmini0898-answerchoice}
  A displayed force is strictly nearer than every alternative
\end{definition}
\begin{proof}
  This is the declared model definition of Is Unique Nearest Displayed Force.
\end{proof}
\begin{lemma}[source To Target Displacements from figure]
  \label{lem:physics:phyx-mini-0898:phyxminiproblems-problemphyxmini0898-sourcetotargetdisplacements-from-figure}
  \lean{PhyXMiniProblems.ProblemPhyXMini0898.sourceToTargetDisplacements_from_figure}
  \uses{def:physics:phyx-mini-0898:phyxminiproblems-problemphyxmini0898-threepointchargeforcesetup, def:physics:phyx-mini-0898:phyxminiproblems-problemphyxmini0898-sourcetotargetdisplacementinmeters, def:physics:phyx-mini-0898:phyxminiproblems-problemphyxmini0898-matchessuppliedthreechargefigure}
  The primary-image coordinates give the vertical 4 cm displacement to the positive source and the 3-4-5 diagonal displacement to the negative source
\end{lemma}
\begin{proof}
  The conclusion follows from the governing relations and figure readouts named in the statement of source To Target Displacements from figure.
\end{proof}
\begin{lemma}[resultant Force On Target eq coulomb Sum]
  \label{lem:physics:phyx-mini-0898:phyxminiproblems-problemphyxmini0898-resultantforceontarget-eq-coulombsum}
  \lean{PhyXMiniProblems.ProblemPhyXMini0898.resultantForceOnTarget_eq_coulombSum}
  \uses{def:physics:phyx-mini-0898:phyxminiproblems-problemphyxmini0898-chargeincoulombs, def:physics:phyx-mini-0898:phyxminiproblems-problemphyxmini0898-forcevectorinnewtons, def:physics:phyx-mini-0898:phyxminiproblems-problemphyxmini0898-sourcecharge, def:physics:phyx-mini-0898:phyxminiproblems-problemphyxmini0898-sourcechargelabel, def:physics:phyx-mini-0898:phyxminiproblems-problemphyxmini0898-targetchargelabel, def:physics:phyx-mini-0898:phyxminiproblems-problemphyxmini0898-threepointchargeforcesetup, def:physics:phyx-mini-0898:phyxminiproblems-problemphyxmini0898-sourcetotargetdisplacementinmeters, def:physics:phyx-mini-0898:phyxminiproblems-problemphyxmini0898-hasphysicalthreechargeparameters, def:physics:phyx-mini-0898:phyxminiproblems-problemphyxmini0898-satisfiespairwisecoulombforcelaw, def:physics:phyx-mini-0898:phyxminiproblems-problemphyxmini0898-satisfieselectrostaticforcesuperposition}
  Coulomb's law and superposition express the resultant in terms of only the independent source data and the governing electromagnetic constant
\end{lemma}
\begin{proof}
  The conclusion follows from the governing relations and figure readouts named in the statement of resultant Force On Target eq coulomb Sum.
\end{proof}
\begin{lemma}[resultant Force Magnitude numerical Bounds]
  \label{lem:physics:phyx-mini-0898:phyxminiproblems-problemphyxmini0898-resultantforcemagnitude-numericalbounds}
  \lean{PhyXMiniProblems.ProblemPhyXMini0898.resultantForceMagnitude_numericalBounds}
  \uses{def:physics:phyx-mini-0898:phyxminiproblems-problemphyxmini0898-forcemagnitudeinnewtons, def:physics:phyx-mini-0898:phyxminiproblems-problemphyxmini0898-threepointchargeforcesetup, def:physics:phyx-mini-0898:phyxminiproblems-problemphyxmini0898-matchessuppliedthreechargefigure, def:physics:phyx-mini-0898:phyxminiproblems-problemphyxmini0898-usesschoolcoulombconstant, def:physics:phyx-mini-0898:phyxminiproblems-problemphyxmini0898-hasphysicalthreechargeparameters, def:physics:phyx-mini-0898:phyxminiproblems-problemphyxmini0898-satisfiespairwisecoulombforcelaw, def:physics:phyx-mini-0898:phyxminiproblems-problemphyxmini0898-satisfieselectrostaticforcesuperposition}
  For the three charge labels and rectangle dimensions in the image, the resultant magnitude is approximately 1.7465e-4 N
\end{lemma}
\begin{proof}
  The conclusion follows from the governing relations and figure readouts named in the statement of resultant Force Magnitude numerical Bounds.
\end{proof}
% --- Archon declaration topology end ---
% --- Archon physics formalization source end ---
